\documentclass[../DoAn.tex]{subfiles}
\begin{document}

Chương 1 giới thiệu tổng quan về đề tài nghiên cứu, bao gồm bối cảnh và sự cần thiết của hệ thống theo dõi vị trí trẻ em theo thời gian thực, mục tiêu và phạm vi đặt ra, định hướng giải pháp kỹ thuật, và cấu trúc tổng thể của báo cáo.

\section{Đặt vấn đề}
\label{section:1.1}

An toàn trẻ em luôn là mối quan tâm hàng đầu của các bậc phụ huynh trong xã hội hiện đại. Theo thống kê của Cục Trẻ em (Bộ Lao động – Thương binh và Xã hội), mỗi năm tại Việt Nam ghi nhận hàng nghìn trường hợp trẻ em bị mất tích, lạc đường hoặc gặp tai nạn bất ngờ ở các khu vực công cộng. Phần lớn các sự việc xảy ra tại các đô thị lớn với mật độ giao thông cao, nơi trẻ có thể dễ dàng bị lạc hoặc rơi vào các tình huống nguy hiểm. Thời gian phát hiện và xử lý chậm trễ trong những tình huống này có thể để lại hậu quả không thể khắc phục.

Bên cạnh đó, nhiều phụ huynh phải để con tự di chuyển đến trường học, trung tâm học thêm hoặc các hoạt động ngoại khóa mà không thể trực tiếp giám sát. Đặc biệt đối với trẻ em trong độ tuổi tiểu học và trung học cơ sở, các em đã có khả năng di chuyển độc lập nhưng chưa đủ nhận thức để bảo vệ bản thân khỏi những nguy cơ tiềm ẩn. Trong thực tế, nhiều phụ huynh chỉ phát hiện con bị lạc hoặc không về đúng giờ khi đã quá trễ để xử lý kịp thời. Nhu cầu theo dõi vị trí của trẻ theo thời gian thực, kết hợp với cơ chế cảnh báo tự động khi trẻ rời khỏi vùng an toàn đã được xác định trước, là bài toán cấp thiết cần giải quyết.

Sự phát triển mạnh mẽ của công nghệ Internet of Things (IoT) và hạ tầng viễn thông di động trong những năm gần đây đã mở ra cơ hội xây dựng các hệ thống theo dõi vị trí nhỏ gọn, chi phí thấp và hoạt động hoàn toàn độc lập \cite{iot_survey}. Nếu được giải quyết hiệu quả, bài toán này không chỉ mang lại sự yên tâm cho hàng triệu phụ huynh mà còn có thể được mở rộng sang các lĩnh vực liên quan như giám sát người cao tuổi mắc chứng mất trí nhớ, quản lý học sinh trong các chuyến dã ngoại tập thể, hoặc theo dõi tài sản di động.

\section{Mục tiêu và phạm vi đề tài}
\label{section:1.2}

Hiện nay đã có một số giải pháp theo dõi vị trí trẻ em trên thị trường. Ứng dụng Life360 \cite{life360} cho phép phụ huynh theo dõi vị trí của điện thoại thông minh mà trẻ mang theo. Giải pháp này hoàn toàn phụ thuộc vào việc trẻ phải có điện thoại đã cài ứng dụng và kết nối Internet thường xuyên — điều không khả thi với trẻ nhỏ hoặc khi điện thoại bị tịch thu tại trường. Apple AirTag và các thiết bị Bluetooth tương tự sử dụng mạng lưới thiết bị lân cận để định vị, nhưng phạm vi hoạt động hạn chế và không hỗ trợ theo dõi liên tục theo thời gian thực. Các GPS tracker thương mại chuyên dụng như Optimus 2.0 hay Bouncie tích hợp SIM di động để truyền vị trí, song thường đi kèm phí đăng ký dịch vụ hàng tháng và thiếu khả năng tùy chỉnh theo nhu cầu cụ thể. Từ phân tích trên, các hạn chế chính của giải pháp hiện tại gồm: phụ thuộc vào điện thoại của trẻ, không theo dõi liên tục thời gian thực, chi phí vận hành cao, và khả năng tùy biến thấp.

Đề tài này đặt mục tiêu xây dựng một hệ thống theo dõi vị trí trẻ em theo thời gian thực hoàn chỉnh, bao gồm ba thành phần chính: (i) thiết bị phần cứng IoT chuyên dụng tích hợp GPS và kết nối di động, hoạt động độc lập không cần điện thoại của trẻ; (ii) hệ thống backend xử lý và lưu trữ dữ liệu vị trí; và (iii) ứng dụng di động dành cho phụ huynh với khả năng xem vị trí thời gian thực, thiết lập vùng an toàn địa lý (geofence), và nhận cảnh báo tức thì khi trẻ vi phạm vùng an toàn.

Phạm vi đề tài tập trung vào xây dựng nguyên mẫu hoàn chỉnh có khả năng vận hành thực tế, bao gồm thiết kế phần cứng, phát triển firmware, backend và ứng dụng di động. Đề tài không bao gồm thiết kế PCB tùy chỉnh hay đóng gói sản phẩm thương mại.

\section{Định hướng giải pháp}
\label{section:1.3}

Từ việc xác định rõ các yêu cầu ở phần \ref{section:1.2}, đề tài lựa chọn hướng tiếp cận xây dựng hệ thống IoT nhúng theo kiến trúc ba tầng: tầng thiết bị, tầng backend và tầng ứng dụng di động.

Về phần cứng, vi điều khiển ESP32 được chọn làm đơn vị xử lý trung tâm nhờ kiến trúc lõi kép (dual-core Xtensa LX6), hiệu năng xử lý cao và hỗ trợ kết nối không dây tích hợp. Module SIM7600 tích hợp đồng thời bộ thu GNSS và modem di động LTE Cat-4 trong một linh kiện duy nhất, cho phép thiết bị kết nối internet qua mạng di động (APN Viettel) và lấy tọa độ GPS mà không cần thêm module ngoài.

Về giao thức truyền thông, MQTT (Message Queuing Telemetry Transport) được chọn để truyền dữ liệu vị trí từ thiết bị lên server. MQTT là giao thức nhẹ theo mô hình publish/subscribe, thiết kế đặc biệt cho thiết bị IoT có tài nguyên hạn chế và kết nối không ổn định, đảm bảo độ trễ thấp và độ tin cậy cao với cơ chế QoS.

Về backend, framework FastAPI (Python) được sử dụng để xây dựng API RESTful và xử lý logic nghiệp vụ bất đồng bộ. Dữ liệu vị trí được lưu trong MongoDB Atlas — cơ sở dữ liệu NoSQL linh hoạt, phù hợp với cấu trúc JSON của payload IoT và hỗ trợ chỉ mục địa lý cho truy vấn lịch sử hành trình.

Về ứng dụng di động, Flutter (Dart) được chọn để xây dựng ứng dụng đa nền tảng Android và iOS từ một codebase duy nhất. Ứng dụng nhận cập nhật vị trí theo thời gian thực từ backend qua WebSocket và hiển thị trên bản đồ OpenStreetMap qua thư viện flutter\_map.

Đóng góp chính của đề tài là một nguyên mẫu hệ thống IoT theo dõi vị trí hoàn chỉnh, hoạt động độc lập không cần điện thoại của trẻ, với độ trễ end-to-end dưới 5 giây và độ chính xác GPS khoảng 3{,}2\,m ở chu kỳ cập nhật 5 giây.

\section{Bố cục đồ án}
\label{section:1.4}

Phần còn lại của báo cáo đồ án tốt nghiệp này được tổ chức như sau.

Chương 2 trình bày kết quả khảo sát hiện trạng các giải pháp theo dõi vị trí trẻ em hiện có, từ đó phân tích và đặc tả các yêu cầu chức năng và phi chức năng của hệ thống cần xây dựng, bao gồm các biểu đồ use case và đặc tả chi tiết các kịch bản sử dụng chính.

Trong Chương 3, đề tài giới thiệu nền tảng lý thuyết và các công nghệ được lựa chọn để triển khai hệ thống, bao gồm vi điều khiển ESP32, module SIM7600, giao thức MQTT, framework FastAPI, cơ sở dữ liệu MongoDB và framework Flutter. Với mỗi công nghệ, chương này phân tích lý do lựa chọn so với các phương án thay thế tương đương.

Chương 4 trình bày toàn bộ quá trình phân tích thiết kế, triển khai và đánh giá hệ thống, bao gồm kiến trúc tổng thể, thiết kế firmware, thiết kế API và cơ sở dữ liệu, giao diện ứng dụng di động, kết quả kiểm thử và mô hình triển khai thực tế trên đám mây.

Chương 5 tập trung trình bày các giải pháp kỹ thuật và đóng góp nổi bật của đề tài, gồm thiết kế tích hợp phần cứng một module, thuật toán geofencing hai chế độ, kiến trúc truyền thông lai MQTT–WebSocket và cơ chế cảnh báo đa kênh song song.

Chương 6 tổng kết các kết quả đạt được, so sánh với mục tiêu ban đầu, chỉ ra các hạn chế còn tồn tại và đề xuất các hướng phát triển tiếp theo cho hệ thống.

\end{document}
